% !TeX root = ../Book1.tex
\begin{figure}[htbp]
	\centering
	\begin{tikzpicture}[x=1.25cm,y=0.9cm,font=\small]
		\fill[GeometryBlue!7] (0,-1.5) rectangle (3,1.5);
		\draw[dashed,GeometryGray] (0,-1.5)--(3,-1.5);
		\draw[dashed,GeometryGray] (0,1.5)--(3,1.5);
		\draw[thick,GeometryBlue] (0,-1.5)--(0,1.5);
		\draw[thick,GeometryBlue] (3,-1.5)--(3,1.5);
		\draw[->,very thick,GeometryBlue] (0,-0.8)--(0,-0.15);
		\draw[->,very thick,GeometryBlue] (3,-0.8)--(3,-0.15);
		\draw[->,thick] (0,0.45)--(3,0.45);
		\fill (0,0.45) circle (1.3pt);
		\fill (3,0.45) circle (1.3pt);
		\node[above] at (1.5,0.45) {\(\ell\)};
		\node[left] at (0,0.45) {\((0,t)\)};
		\node[right] at (3,0.45) {\((\ell,t)\)};
		\draw[->,GeometryGray] (0,1.85)--(3,1.85);
		\node[above] at (1.5,1.85) {\(T(s,t)=(s+\ell,t)\)};
		\node[below] at (0,-1.5) {\(s=0\)};
		\node[below] at (3,-1.5) {\(s=\ell\)};
	\end{tikzpicture}
	\caption[平圆柱的基本带与平移周期]{平圆柱由基本带的两条竖边按同高点粘合得到，两边箭头同向。横线段投影为一次闭路，在覆盖平面上的端点相差平移 \(\ell\)。上下虚线只表示绘图截断，实际基本带沿 \(t\) 方向无限延伸。}
	\label{g1:fig:euclidean-cylinder-strip}
\end{figure}
